% 教学任务地图；宽约108mm，采用项目语义配色。
\begin{tikzpicture}[
  font=\figurefont\small,
  task/.style={draw=BookRule,fill=BookPaper,minimum width=30mm,
    minimum height=9mm,align=center,inner sep=2mm},
  link/.style={-{Stealth[length=2mm]},draw=BookMuted,line width=0.8pt},
  note/.style={anchor=west,text=BookMuted,font=\figurefont\footnotesize}
]
  \node[task,fill=FigureInput!10,draw=FigureInput,minimum width=20mm] (input) at (0,0) {街景观测};
  \node[task] (recognition) at (4,2.6) {类别、属性};
  \node[task] (region) at (4,1.3) {框、掩码、轨迹};
  \node[task] (measure) at (4,0) {数量、距离、排序};
  \node[task] (restore) at (4,-1.3) {复原、压缩结果};
  \node[task] (create) at (4,-2.6) {编辑、描述、配声};
  \draw[draw=BookMuted,line width=0.8pt] (input.east) -- (1.8,0);
  \draw[draw=BookMuted,line width=0.8pt] (1.8,-2.6) -- (1.8,2.6);
  \foreach \name/\y in {recognition/2.6,region/1.3,measure/0,restore/-1.3,create/-2.6}{
    \draw[link] (1.8,\y) -- (\name.west);
  }
  \node[note] at (5.8,2.6) {对象与语义};
  \node[note] at (5.8,1.3) {范围与身份};
  \node[note] at (5.8,0) {尺度与相关性};
  \node[note] at (5.8,-1.3) {保真与比特};
  \node[note] at (5.8,-2.6) {条件与对应};
\end{tikzpicture}
